% GENERATED FILE. Edit canonical lessons, then run npm run generate:books.
% canonical-source-hash: fnv1a64:c744bf80a9e0c0c8
% canonical-lessons: JA-C64-kona, JA-C64-goma, JA-C64-uni, JA-C64-iwashi, JA-C64-aji

\chapter{Powder, Sesame, Sea urchin, Sardine, Flavour}
\label{ch:ja-powder-sesame-sea-urchin-sardine-flavour}
\hlchaptermodality{64}

\begin{chapteropening}
\textbf{By the end of this chapter:} \emph{I can say powder, flour, sesame, a sea urchin, a sardine and a flavour, a taste.}

Complete the last lesson of chapter 64: I can say powder, flour, sesame, a sea urchin, a sardine and a flavour, a taste.
\end{chapteropening}

\section[\texorpdfstring{\ja{こな} (kona)}{こな (kona)}]{\ja{こな} (kona) — powder, flour}

\label{lesson:JA-C64-kona}



\begin{quote}
Before the new one: say the Japanese for east, then the Japanese for south.
\end{quote}

\subsection*{You'll want to know: \ja{こな}}
\textbf{\ja{こな}} — \emph{kona} — \textquotedblleft{}powder, flour\textquotedblright{}.

\begin{cousinweb}
\textbf{\ja{こな}} (\emph{kona}) is a native Japanese word, inherited rather than borrowed from Chinese, and its older origin is not settled. Its kanji was attached to it later as a kun'yomi, a native reading given to a Chinese character (JMdict, KANJIDIC2).
\end{cousinweb}

\begin{grammarlens}[title={What you've built}]
One more word for this chapter.
\end{grammarlens}

\subsection*{Your turn}
\begin{itemize}
\raggedright
  \item \emph{Say it:} \emph{kona}
  \item \emph{Say it:} \emph{kona}, once more
  \item \emph{Say it:} say \emph{kona}
  \item \emph{Recall:} say the Japanese for east, then the Japanese for south, then say \emph{kona} again
\end{itemize}

\begin{tcolorbox}[breakable,colback=teal!4,colframe=teal!35!black,title={Before you move on}]
What does \ja{こな} mean? (\textquotedblleft{}Powder, flour\textquotedblright{}.) Say it once more.
\end{tcolorbox}
\section[\texorpdfstring{\ja{ごま} (goma)}{ごま (goma)}]{\ja{ごま} (goma) — sesame}

\label{lesson:JA-C64-goma}



\begin{quote}
Before the new one: say the Japanese for south, then the Japanese for powder.
\end{quote}

\subsection*{You'll want to know: \ja{ごま}}
\textbf{\ja{ごま}} — \emph{goma} — \textquotedblleft{}sesame\textquotedblright{}.

\begin{cousinweb}
\textbf{\ja{ごま}} (\emph{goma}) is Sino-Japanese (kango): two Chinese characters read with their on'yomi, readings borrowed from Chinese. They are \emph{go}, \textquotedblleft{}foreign\textquotedblright{} (Middle Chinese \emph{hu}) and \emph{ma}, \textquotedblleft{}hemp\textquotedblright{} (\emph{mae}) (Unihan, KANJIDIC2; Middle Chinese after Baxter–Sagart). JMdict marks it as usually written in kana.
\end{cousinweb}

\begin{grammarlens}[title={What you've built}]
One more word for this chapter.
\end{grammarlens}

\subsection*{Your turn}
\begin{itemize}
\raggedright
  \item \emph{Say it:} \emph{goma}
  \item \emph{Say it:} \emph{goma}, once more
  \item \emph{Say it:} say \emph{goma}
  \item \emph{Recall:} say the Japanese for south, then the Japanese for powder, then say \emph{goma} again
\end{itemize}

\begin{tcolorbox}[breakable,colback=teal!4,colframe=teal!35!black,title={Before you move on}]
What does \ja{ごま} mean? (\textquotedblleft{}Sesame\textquotedblright{}.) Say it once more.
\end{tcolorbox}
\section[\texorpdfstring{\ja{うに} (uni)}{うに (uni)}]{\ja{うに} (uni) — a sea urchin}

\label{lesson:JA-C64-uni}



\begin{quote}
Before the new one: say the Japanese for powder, then the Japanese for sesame.
\end{quote}

\subsection*{You'll want to know: \ja{うに}}
\textbf{\ja{うに}} — \emph{uni} — \textquotedblleft{}a sea urchin\textquotedblright{}.

\begin{cousinweb}
\textbf{\ja{うに}} (\emph{uni}) is a native word, inherited rather than borrowed. JMdict writes it with characters for \textquotedblleft{}sea\textquotedblright{} and \textquotedblleft{}gall\textquotedblright{}, read as a whole, and marks it as usually written in kana. Its older origin is not settled.
\end{cousinweb}

\begin{grammarlens}[title={What you've built}]
One more word for this chapter.
\end{grammarlens}

\subsection*{Your turn}
\begin{itemize}
\raggedright
  \item \emph{Say it:} \emph{uni}
  \item \emph{Say it:} \emph{uni}, once more
  \item \emph{Say it:} say \emph{uni}
  \item \emph{Recall:} say the Japanese for powder, then the Japanese for sesame, then say \emph{uni} again
\end{itemize}

\begin{tcolorbox}[breakable,colback=teal!4,colframe=teal!35!black,title={Before you move on}]
What does \ja{うに} mean? (\textquotedblleft{}A sea urchin\textquotedblright{}.) Say it once more.
\end{tcolorbox}
\section[\texorpdfstring{\ja{いわし} (iwashi)}{いわし (iwashi)}]{\ja{いわし} (iwashi) — a sardine}

\label{lesson:JA-C64-iwashi}



\begin{quote}
Before the new one: say the Japanese for sesame, then the Japanese for a sea urchin.
\end{quote}

\subsection*{You'll want to know: \ja{いわし}}
\textbf{\ja{いわし}} — \emph{iwashi} — \textquotedblleft{}a sardine\textquotedblright{}.

\begin{cousinweb}
\textbf{\ja{いわし}} (\emph{iwashi}) is a native Japanese word, inherited rather than borrowed from Chinese, and its older origin is not settled. It is usually written in kana; its kanji, when used, is a kun'yomi attached later (JMdict, KANJIDIC2).
\end{cousinweb}

\begin{grammarlens}[title={What you've built}]
One more word for this chapter.
\end{grammarlens}

\subsection*{Your turn}
\begin{itemize}
\raggedright
  \item \emph{Say it:} \emph{iwashi}
  \item \emph{Say it:} \emph{iwashi}, once more
  \item \emph{Say it:} say \emph{iwashi}
  \item \emph{Recall:} say the Japanese for sesame, then the Japanese for a sea urchin, then say \emph{iwashi} again
\end{itemize}

\begin{tcolorbox}[breakable,colback=teal!4,colframe=teal!35!black,title={Before you move on}]
What does \ja{いわし} mean? (\textquotedblleft{}A sardine\textquotedblright{}.) Say it once more.
\end{tcolorbox}
\section[\texorpdfstring{\ja{あじ} (aji)}{あじ (aji)}]{\ja{あじ} (aji) — a flavour, a taste}

\label{lesson:JA-C64-aji}



\begin{quote}
Before the new one: say the Japanese for a sea urchin, then the Japanese for a sardine.
\end{quote}

\subsection*{You'll want to know: \ja{あじ}}
\textbf{\ja{あじ}} — \emph{aji} — \textquotedblleft{}a flavour, a taste\textquotedblright{}.

\begin{cousinweb}
\textbf{\ja{あじ}} (\emph{aji}) is a native Japanese word, inherited rather than borrowed from Chinese, and its older origin is not settled. Its kanji was attached to it later as a kun'yomi, a native reading given to a Chinese character (JMdict, KANJIDIC2).
\end{cousinweb}

\begin{grammarlens}[title={What you've built}]
One more word for this chapter.
\end{grammarlens}

\subsection*{Your turn}
\begin{itemize}
\raggedright
  \item \emph{Say it:} \emph{aji}
  \item \emph{Say it:} \emph{aji}, once more
  \item \emph{Say it:} say \emph{aji}
  \item \emph{Recall:} say the Japanese for a sea urchin, then the Japanese for a sardine, then say \emph{aji} again
\end{itemize}

\begin{tcolorbox}[breakable,colback=teal!4,colframe=teal!35!black,title={Before you move on}]
What does \ja{あじ} mean? (\textquotedblleft{}A flavour, a taste\textquotedblright{}.) Say it once more.
\end{tcolorbox}
